\section{AI for Math: razonamiento formal y LLM Agent}

\subsection{Limitaciones de los métodos puramente neuronales}

En años recientes los LLM han tenido un desempeño sobresaliente en competencias de matemáticas (por ejemplo, la mejora notable de O1 en el AMC 2024, el desempeño de nivel medalla de plata de AlphaProof en la IMO), pero los métodos puramente neuronales tienen debilidades fundamentales:

\begin{warningbox}{Dos problemas fundamentales de los modelos de lenguaje puros para hacer matemáticas}
\begin{enumerate}
    \item \textbf{Escasez de datos}: el razonamiento matemático riguroso requiere una gran cantidad de trayectorias de razonamiento de alta calidad o funciones de recompensa de alta calidad; más allá del nivel de competencias o de preparatoria, es difícil obtenerlas solo a partir de datos humanos
    \item \textbf{El razonamiento en lenguaje natural es difícil de verificar}: la corrección del razonamiento en lenguaje no puede comprobarse formalmente, lo cual es especialmente riesgoso en sistemas reales donde los casos límite son críticos
\end{enumerate}
\end{warningbox}

\subsection{Representación formal: una alternativa}

\begin{importantbox}{Idea central de los métodos formales}
\begin{enumerate}
    \item \textbf{Autoformalización} (Autoformalization): convertir enunciados matemáticos informales al lenguaje formal de sistemas como Lean, Coq o Isabelle
    \item \textbf{Demostrador neuronal} (Neural Prover): busca secuencias de estrategias de demostración (tactics) dentro de un entorno formal
    \item \textbf{Verificación}: una demostración formal puede comprobarse automáticamente mediante un asistente de demostración; esto resuelve el problema de la verificación
    \item \textbf{Datos sintéticos}: se pueden generar en gran cantidad demostraciones candidatas y obtener de inmediato retroalimentación sobre su corrección, lo que convierte el problema de demostración en una tarea similar a un juego
\end{enumerate}
\end{importantbox}

\subsubsection{Cómo funciona el asistente de demostración Lean}

Lean es una máquina de estados:
\begin{itemize}
    \item \textbf{Estado} (State): los objetivos (goals) que se necesitan demostrar en el momento actual y las hipótesis del contexto
    \item \textbf{Acción} (Tactic): reglas de simplificación que, al aplicarse, cambian el estado de la demostración
    \item \textbf{Objetivo}: encontrar una secuencia de tactics que resuelva todos los goals y llegue al estado QED
\end{itemize}

Por ejemplo, para demostrar «si $x$ es par, entonces $x^2$ es par»: el estado inicial contiene el objetivo $x^2 \bmod 2 = 0$; mediante tactics como \texttt{intro} y \texttt{rw} se simplifica paso a paso hasta llegar a $0 = 0$, es decir, QED.

\subsection{Copra: demostración de teoremas basada en LLM Agent}

\begin{importantbox}{Arquitectura del sistema Copra}
Copra es un Agent de demostración de teoremas basado en aprendizaje en contexto (In-Context Learning):
\begin{enumerate}
    \item Usa un LLM de vanguardia (como GPT-4) para predecir la siguiente tactic
    \item Proporciona el estado de la demostración y el historial de búsqueda al LLM como prompt
    \item Ejecuta la tactic predicha y obtiene el nuevo estado o retroalimentación de error
    \item Incorpora la información del error de vuelta al prompt para que el LLM corrija su estrategia
    \item En el nivel externo usa búsqueda en profundidad (DFS) para organizar todo el proceso de demostración
\end{enumerate}
\end{importantbox}

Ventajas centrales de Copra:
\begin{itemize}
    \item \textbf{No requiere entrenamiento adicional}: aprovecha directamente las capacidades de un LLM preentrenado
    \item \textbf{Se beneficia de inmediato de los avances de los LLM}: la actualización de GPT-4 a O1/O3 mejora el desempeño directamente
    \item \textbf{Combina el razonamiento en lenguaje natural con el razonamiento en lenguaje formal}
    \item \textbf{No requiere un corpus de entrenamiento}: es aplicable a problemas de investigación completamente nuevos
    \item \textbf{La retroalimentación del asistente de demostración provee un anclaje} (grounding): evita las alucinaciones
\end{itemize}

\subsection{Demostración jerárquica: el poder de la abstracción}

\begin{importantbox}{Descomposición jerárquica de demostraciones impulsada por LLM}
Para teoremas complejos (como los problemas de la IMO), Copra puede extenderse a una resolución jerárquica:
\begin{enumerate}
    \item Se le pide al LLM que genere un plan de demostración informal (informal proof plan)
    \item A partir del plan, se descompone el teorema en la formulación formal de varios subobjetivos (sub-goals)
    \item Se usa Copra para demostrar los subobjetivos uno por uno
    \item Se colocan los subobjetivos ya demostrados en el contexto y se demuestra el teorema completo
\end{enumerate}
Idea clave: el LLM no solo puede predecir pasos de demostración de bajo nivel, sino también realizar \textbf{razonamiento abstracto de alto nivel}; todo esto se logra mediante ingeniería de prompt, sin necesidad de entrenamiento adicional.
\end{importantbox}

Tomando como ejemplo el problema de la IMO «demostrar que $\frac{21n+4}{14n+3}$ es irreducible para todo número natural $n$»: O3 descompone el problema en tres lemas auxiliares sobre el GCD; después de que Copra los demuestra por separado, se combinan para demostrar el teorema completo.

\subsection{Aplicación: verificación formal de compiladores}

\begin{knowledgebox}{El valor práctico de la verificación formal}
La verificación formal puede demostrar matemáticamente la corrección, la seguridad y el desempeño de un sistema. El Departamento de Defensa de Estados Unidos llegó a desplegar en drones un núcleo Linux verificado formalmente (seL4), y un equipo rojo no pudo vulnerarlo en seis semanas.

Históricamente la verificación formal ha sido extremadamente costosa, pero el auge de AI for Math está \textbf{cambiando de manera fundamental la ecuación de costo-beneficio}.
\end{knowledgebox}

La clase mostró un caso sencillo de verificación de un compilador de expresiones aritméticas (en el lenguaje Coq):
\begin{itemize}
    \item Lenguaje fuente: expresiones aritméticas definidas de manera recursiva (suma, multiplicación)
    \item Lenguaje objetivo: una lista de instrucciones basada en una máquina de pila
    \item Teorema de corrección del compilador: el resultado de ejecutar el código compilado es igual al resultado de evaluar el lenguaje fuente
    \item Copra no puede demostrarlo de una sola vez, pero mediante el LLM se inventa un lema auxiliar (la corrección de la compilación de una sola instrucción) y, al combinarlos, se completa la demostración completa
\end{itemize}

\subsection{Resumen de la sección}

\begin{itemize}
    \item Los LLM Agent para el descubrimiento matemático son muy prometedores; los métodos de aprendizaje en contexto ya han mostrado capacidades sólidas
    \item La retroalimentación del asistente de demostración es fundamental para el anclaje (grounding)
    \item Un diseño de Agent avanzado puede combinar la formalización, la planificación jerárquica y la generación de demostraciones de bajo nivel
    \item Por ahora todavía se necesita una arquitectura de Agent; los métodos puramente de entrenamiento aún no alcanzan un nivel suficiente
\end{itemize}
